\section{Longest Common Substring}
\label{sec:dp-longest-common-substring}

\problemstatement{Given two strings $s_1$ and $s_2$, find the length of
their longest common \textbf{substring}: a longest run of
\textbf{contiguous} characters that appears, unbroken, in both strings.}

\begin{patternbox}
The keyword shift from \emph{subsequence} (Section~\ref{sec:dp-lcs}) to
\emph{substring} -- contiguity now required -- breaks the mismatch case
of the LCS recurrence entirely: a substring cannot ``skip'' a
mismatched character and continue, so a mismatch must \textbf{reset} the
current run to length $0$, rather than falling back to the best of two
neighboring sub-answers.
\end{patternbox}

\subsection*{Deriving the recurrence}

Let $f(i,j)$ be the length of the common substring \textbf{ending
exactly at} $s_1[i-1]$ and $s_2[j-1]$ (not the best answer over all
prefixes, as in LCS -- specifically the run ending right here):
\[
  f(i,j) =
  \begin{cases}
    1 + f(i{-}1,j{-}1) & \text{if } s_1[i{-}1] = s_2[j{-}1], \\
    0 & \text{otherwise.}
  \end{cases}
\]
Because the answer is not necessarily at $f(n,m)$ (the longest common
run could end anywhere in either string), we additionally track a
running global maximum over every $f(i,j)$ computed.

\subsection*{Approach 1: Recursion}

\begin{ideabox}[Why a Direct Recursive Translation Is Awkward Here]
Unlike every previous problem in this book, $f(i,j)$ alone does not
directly answer the original question -- the answer is
$\max_{i,j} f(i,j)$, a quantity that is naturally accumulated during a
forward (tabulation-style) sweep, not easily expressed as a single
top-down recursive call. For this reason, Longest Common Substring is
typically presented starting from tabulation, with the recursive/memoized
forms understood as computing the same per-cell $f(i,j)$ values, just
collected via top-down calls instead of a bottom-up loop.
\end{ideabox}

\subsection*{Approach 2: Memoization}

\begin{lstlisting}
#include <bits/stdc++.h>
using namespace std;

int f(int i, int j, string& s1, string& s2, vector<vector<int>>& memo, int& best) {
    if (i == 0 || j == 0) return 0;
    if (memo[i][j] != -1) return memo[i][j];

    int result = 0;
    if (s1[i - 1] == s2[j - 1]) {
        result = 1 + f(i - 1, j - 1, s1, s2, memo, best);
    }
    best = max(best, result);
    return memo[i][j] = result;
}

int longestCommonSubstringMemo(string& s1, string& s2) {
    int n = s1.size(), m = s2.size();
    vector<vector<int>> memo(n + 1, vector<int>(m + 1, -1));
    int best = 0;

    for (int i = 1; i <= n; i++) {
        for (int j = 1; j <= m; j++) {
            f(i, j, s1, s2, memo, best);
        }
    }
    return best;
}

int main() {
    string s1 = "abcde", s2 = "abfce";
    cout << longestCommonSubstringMemo(s1, s2) << endl; // 2
    return 0;
}
\end{lstlisting}

\begin{complexitybox}[Approach 2 Complexity]
\textbf{Time.} $O(n \cdot m)$.
\textbf{Space.} $O(n \cdot m)$ memo.
\end{complexitybox}

\subsection*{Approach 3: Tabulation}

\begin{lstlisting}
#include <bits/stdc++.h>
using namespace std;

int longestCommonSubstringTabulation(string& s1, string& s2) {
    int n = s1.size(), m = s2.size();
    vector<vector<int>> dp(n + 1, vector<int>(m + 1, 0));
    int best = 0;

    for (int i = 1; i <= n; i++) {
        for (int j = 1; j <= m; j++) {
            if (s1[i - 1] == s2[j - 1]) {
                dp[i][j] = 1 + dp[i - 1][j - 1];
            } else {
                dp[i][j] = 0;
            }
            best = max(best, dp[i][j]);
        }
    }
    return best;
}

int main() {
    string s1 = "abcde", s2 = "abfce";
    cout << longestCommonSubstringTabulation(s1, s2) << endl; // 2
    return 0;
}
\end{lstlisting}

\subsection*{Dry run of the tabulation table}

\dryrun{Take $s_1=\texttt{"abcde"}$, $s_2=\texttt{"abfce"}$.}

\begin{center}
\begin{tabular}{c|c|c|c|c|c|c}
 & \texttt{""} & \texttt{a} & \texttt{b} & \texttt{f} & \texttt{c} & \texttt{e} \\ \hline
\texttt{""} & 0 & 0 & 0 & 0 & 0 & 0 \\ \hline
\texttt{a} & 0 & 1 & 0 & 0 & 0 & 0 \\ \hline
\texttt{b} & 0 & 0 & 2 & 0 & 0 & 0 \\ \hline
\texttt{c} & 0 & 0 & 0 & 0 & 1 & 0 \\ \hline
\texttt{d} & 0 & 0 & 0 & 0 & 0 & 0 \\ \hline
\texttt{e} & 0 & 0 & 0 & 0 & 0 & 1
\end{tabular}
\end{center}

The running maximum reaches $2$ at cell (\texttt{b}, \texttt{b}) --
the common substring \texttt{"ab"} -- and never exceeds it elsewhere
(every other match is isolated, contributing only length $1$ runs).

\begin{complexitybox}[Approach 3 Complexity]
\textbf{Time.} $O(n \cdot m)$.
\textbf{Space.} $O(n \cdot m)$ table, no recursion stack.
\end{complexitybox}

\subsection*{Approach 4: Space Optimization}

\begin{lstlisting}
#include <bits/stdc++.h>
using namespace std;

int longestCommonSubstringSpaceOptimized(string& s1, string& s2) {
    int n = s1.size(), m = s2.size();
    vector<int> prev(m + 1, 0), current(m + 1, 0);
    int best = 0;

    for (int i = 1; i <= n; i++) {
        for (int j = 1; j <= m; j++) {
            if (s1[i - 1] == s2[j - 1]) {
                current[j] = 1 + prev[j - 1];
            } else {
                current[j] = 0;
            }
            best = max(best, current[j]);
        }
        prev = current;
    }
    return best;
}

int main() {
    string s1 = "abcde", s2 = "abfce";
    cout << longestCommonSubstringSpaceOptimized(s1, s2) << endl; // 2
    return 0;
}
\end{lstlisting}

\begin{complexitybox}[Approach 4 Complexity]
\textbf{Time.} $O(n \cdot m)$.
\textbf{Space.} $O(m)$ for two rolling rows.
\end{complexitybox}

\begin{takeawaybox}
``Subsequence'' (gaps allowed) and ``substring'' (no gaps allowed) are
the single most important keyword distinction in string DP, and they
produce genuinely different recurrences -- LCS falls back to a
neighboring best-of-two on a mismatch, while Longest Common Substring
resets to zero, because a broken run cannot be patched back together.
Always check this one word before assuming a string problem is ``LCS in
disguise.''
\end{takeawaybox}
